\section{Introduction}

The study of polynomial equations is a basic question
of mathematics.  In this paper we study a problem we
call \dproblem{Max-$d$-Eq}
where we are given a system of $m$ equations of degree
$d$ in $n$ variables over $\GF[2]$.  As we consider the
case of $d$ constant, all polynomials are given in
the dense representation, \ie, by a list of monomials.
Many problems can be coded as polynomial equations and in particular
it is easy to code \dproblem{3-Sat} as equations of degree 3
and thus determining whether we can simultaneously
satisfy all equations is \cclass{NP}-complete.   It is hence
natural to study
the question of satisfying the %
maximum
number of
equations and our interest turns to approximation
algorithms.  We say that an algorithm is a $C$-approximation
algorithm if it always returns a solution which satisfies
at least $C \cdot \mbox{OPT}$ equations where
$\mbox{OPT}$ is the number of equations satisfied
by the optimal solution.  
